
\medskip

%\emph{Cet exercice est un questionnaire à choix multiples.\\Pour chacune des questions suivantes, une seule des quatre réponses proposées est exacte.\\Pour répondre, indiquer sur la copie le numéro de la question et la lettre de la réponse choisie.\\Aucune justification n'est demandée.\\Une réponse fausse, une absence de réponse, ou une réponse multiple, ne rapporte
%ni n'enlève de point.}
%
%\medskip

\begin{enumerate}
\item On considère la fonction $f$ définie sur $\R$ par $f(x) = 2x\e^x$.

%Le nombre de solutions sur $\R$ de l'équation $f(x) = - \dfrac{73}{100}$ est égal à : 
La fonction $f$ produit de fonctions dérivables sur $\R$ est dérivable et pour tout réel : $f'(x) = 2\e^x + 2x\e^x = 2\e^x(1 + x)$. Cette dérivée est du signe de $1 + x$.

On a donc $f'(x) < 0 \iff x < - 1$ : la fonction est décroissante sur $]- \infty~;~- 1]$.

On sait que $\displaystyle\lim_{x \to -\infty} x\e^x = 0$, donc $\displaystyle\lim_{x \to -\infty}f(x) = 0$ et $f(-1) = 2 \times (- 1) \times \e{-1} = - 2\e^{-1} \approx - 0,74$.

De même $f'(x) >  \iff x > - 1$ : la fonction est croissante sur $[- 1~;~+ \infty[$.

Or  $\displaystyle\lim_{x \to +\infty}f(x) = + \infty$.

Sur l'intervalle $[- 1~;~+ \infty[$, la fonction $f$ croit de $f(-1) < -0,73$ à plus l'infini.

La fonction $f$ étant continue car dérivable et strictement croissante prend donc une seule fois la valeur $-0,73$ d'après le théorème des valeurs intermédiaires. Réponse \textbf{c.}
%\begin{center}
%\begin{tabularx}{\linewidth}{*{4}{X}}
%\textbf{a.~~} 0 &\textbf{b.~~} 1 &\textbf{c.~~}2&\textbf{d.~~}une infinité.
%\end{tabularx}
%\end{center}

\item %On considère la fonction $g$ définie sur $\R$ par :
\[g(x) = \dfrac{x+ 1}{\e^x}.\]

%La limite de la fonction $g$ en $- \infty$ est égale à : 
%
%\begin{center}
%\begin{tabularx}{\linewidth}{*{4}{X}}
%\textbf{a.~~} $-\infty$&\textbf{b.~~} $+\infty$ &\textbf{c.~~} 0&\textbf{d}elle n'existe pas.
%\end{tabularx}
%\end{center}
On a $g(x) = (x + 1) \times \dfrac{1}{\e^x}$. Or :

$\displaystyle\lim_{x \to -\infty}x + 1= - \infty$ et 

$\displaystyle\lim_{x \to -\infty}\dfrac{1}{\e^x} = + \infty$, donc par produit de limites : $\displaystyle\lim_{x \to -\infty}g(x) = - \infty$. Réponse \textbf{a.}

\item %On considère la fonction $h$ définie sur $\R$ par:
\[h(x) = (4x - 16)\e^{2x}.\]

%On note $\mathcal{C}_h$ la courbe représentative de $h$ dans un repère orthogonal.

%On peut affirmer que:

%\begin{center}
%\begin{tabularx}{\linewidth}{*{2}{X}}
%\textbf{a.~~} $h$ est convexe sur $\R$.&\textbf{b.~~} $\mathcal{C}_h$ possède un point d'inflexion en $x = 3$.\\
%\textbf{c.~~} $h$ est concave sur $\R$.&\textbf{d.~~} $\mathcal{C}_h$ possède un point d'inflexion 
%
%en $x = 3,5$.
%\end{tabularx}
%\end{center}
La fonction $h$ produit de fonctions dérivables sur $\R$ est dérivable et pour tout réel :

$h'(x) = 4\e^{2x} + 2(4x - 16)\e^{2x} = \e^{2x}(4 + 8x - 32) = \e^{2x}(8x - 28)$.

$h''(x) = 8\e^{2x} + 2(8x - 28)\e^{2x} = \e^{2x}(8 + 16x  - 56) = \e^{2x}(16x - 48) = 16\e^{2x}(x - 3)$.

Comme $4\e^{2x} > 0$ quel que soit le réel $x$, le signe de $h''(x)$ est celui de $x - 3$.

Donc la dérivée seconde s'annule en $x = 3$ en changeant de signe, donc le point de coordonnées $\left(3~;~-4\text{e}^{6}\right)$ est le seul point d'inflexion de la courbe $\mathcal{C}_h$. Réponse \textbf{b.}

\item %On considère la fonction $k$ définie sur l'intervalle $]0~;~ +\infty[$ par : 
\[k(x) = 3 \ln (x) - x.\]

%On note $\mathcal{C}$ la courbe représentative de la fonction $k$ dans un repère orthonormé. 

%On note $T$ la tangente à la courbe $\mathcal{C}$ au point d'abscisse $x = \text{e}$.

Une équation de $T$ est : $M(x~;~y) \in (T) \iff y - k(\e) = f'(\e)(x- \e)$.

Avec $k(\e) = 3\ln \e - \e = 3 - \e$ et 

$k'(x) = \dfrac 3x - 1$, d'où $k'(\e) = \dfrac{3}{\e} - 1 = \dfrac{3 - \e}{\e}$, on obtient donc :

$M(x~;~y) \in (T) \iff y - (3 - \e) = \dfrac{3 - \e}{\e}(x- \e) \iff y  = \dfrac{3 - \e}{\e}x - 3 + \e  + 3 - \e \iff y = \dfrac{3 - \e}{\e}x$ (cette droite contient l'origine). Réponse \textbf{b.}

%\begin{center}
%\begin{tabularx}{\linewidth}{*{2}{X}}
%\textbf{a.~~} $y = (3 - \text{e})x$&\textbf{b.~~} $y = \left(\dfrac{3 - \text{e}}{\text{e}}\right)x$\\
%\textbf{c.~~} $y = \left(\dfrac{3}{\text{e} }- 1\right)x + 1$&\textbf{d.~~}$y = (\text{e} - 1)x + 1$
%\end{tabularx}
%\end{center}
\item %On considère l'équation 
$[\ln (x)]^2 + 10 \ln (x) + 21 = 0$, avec $x \in ]0~;~+\infty[$.

%Le nombre de solutions de cette équation est égal à :
%\begin{center}
%\begin{tabularx}{\linewidth}{*{4}{X}}
%\textbf{a.~~} 0&\textbf{b.~~} 1&\textbf{c.~~} 2&\textbf{d.~~} une infinité.
%\end{tabularx}
%\end{center}
En posant $X = \ln x$, l'équation devient : $X^2 + 10X + 21 = 0$.

Avec $\Delta = 10^2 - 4\times 21 = 100 - 84 = 16 = 4^2$, les solutions sont :

$\left\{\begin{array}{l c l}
X_1 &=&\dfrac{-10 - 4}{2}\\
X_2&=&\dfrac{-10 + 4}{2}
\end{array}\right. \iff 
\left\{\begin{array}{l c l}
X_1 &=&- 7\\
X_2&=&- 3
\end{array}\right. \iff \left\{\begin{array}{l c l}
\ln x_1 &=&-7\\
\ln x_2&=&-3
\end{array}\right. \iff 
\left\{\begin{array}{l c l}
x_1 &=&\e^{-7}\\
x_2&=&\e^{-3}
\end{array}\right.$

Cette équation a donc deux solutions. Réponse \textbf{c.}
\end{enumerate}


